% Options for packages loaded elsewhere
\PassOptionsToPackage{unicode}{hyperref}
\PassOptionsToPackage{hyphens}{url}
\documentclass[
  10pt,
  letterpaper,
]{article}
\usepackage{xcolor}
\usepackage[margin=23mm]{geometry}
\usepackage{amsmath,amssymb}
\setcounter{secnumdepth}{5}
\usepackage{iftex}
\ifPDFTeX
  \usepackage[T1]{fontenc}
  \usepackage[utf8]{inputenc}
  \usepackage{textcomp} % provide euro and other symbols
\else % if luatex or xetex
  \usepackage{unicode-math} % this also loads fontspec
  \defaultfontfeatures{Scale=MatchLowercase}
  \defaultfontfeatures[\rmfamily]{Ligatures=TeX,Scale=1}
\fi
\usepackage{lmodern}
\ifPDFTeX\else
  % xetex/luatex font selection
\fi
% Use upquote if available, for straight quotes in verbatim environments
\IfFileExists{upquote.sty}{\usepackage{upquote}}{}
\IfFileExists{microtype.sty}{% use microtype if available
  \usepackage[]{microtype}
  \UseMicrotypeSet[protrusion]{basicmath} % disable protrusion for tt fonts
}{}
\makeatletter
\@ifundefined{KOMAClassName}{% if non-KOMA class
  \IfFileExists{parskip.sty}{%
    \usepackage{parskip}
  }{% else
    \setlength{\parindent}{0pt}
    \setlength{\parskip}{6pt plus 2pt minus 1pt}}
}{% if KOMA class
  \KOMAoptions{parskip=half}}
\makeatother
\setlength{\emergencystretch}{3em} % prevent overfull lines
\providecommand{\tightlist}{%
  \setlength{\itemsep}{0pt}\setlength{\parskip}{0pt}}
\usepackage[]{natbib}
\bibliographystyle{plainnat}
\usepackage{bookmark}
\IfFileExists{xurl.sty}{\usepackage{xurl}}{} % add URL line breaks if available
\urlstyle{same}
\hypersetup{
  pdftitle={An Average Is Not a Guarantee: Making Assumptions Visible},
  pdfauthor={Super Writer worked example},
  hidelinks,
  pdfcreator={LaTeX via pandoc}}

\title{An Average Is Not a Guarantee: Making Assumptions Visible}
\author{Super Writer worked example}
\date{September 2026}

\begin{document}
\maketitle
\begin{abstract}
A bound on an aggregate is easy to overstate when its assumptions
disappear between a proof and an abstract. This educational note
restates the elementary range bound for a convex combination, proves it
directly, and gives a signed- weight counterexample. Nonnegative weights
that sum to one keep a weighted average within the range of its inputs.
The result does not establish that every input is accurate, that the
average is unbiased, or that a predictive model generalizes. The purpose
is to demonstrate assumption-aware theoretical writing, not to present a
new theorem or a conference submission.
\end{abstract}

\section{The Writing Problem}\label{the-writing-problem}

The sentence ``aggregation guarantees reliable predictions'' mixes a
mathematical property with an empirical promise. Even when an aggregate
satisfies a range bound, nothing in that bound establishes that the
inputs are close to an unknown target. The proof and the headline need
to answer the same question.

We use a convex combination, a standard object in convex analysis
\citep{boyd2004}, to separate the two statements. The argument below is
written independently for this example. It is elementary and carries no
novelty claim.

\section{Setting and Proposition}\label{setting-and-proposition}

Let \(y_1,\ldots,y_n\) be real numbers, with finite \(n\geq 1\). Let
\(w_i\geq 0\) for every \(i\), and suppose \(\sum_{i=1}^n w_i=1\).
Define the aggregate

\[
\bar y=\sum_{i=1}^n w_i y_i.
\]

\textbf{Proposition.} Under these assumptions,

\[
\min_i y_i\;\leq\;\bar y\;\leq\;\max_i y_i.
\]

The statement is deterministic. It does not assume a training
distribution, sampling independence, or a probabilistic model.
Conversely, it does not provide any conclusion about those quantities.

\section{Proof and Assumption Use}\label{proof-and-assumption-use}

Write \(m=\min_i y_i\) and \(M=\max_i y_i\). For every index,
\(m\leq y_i\leq M\). Since \(w_i\) is nonnegative, multiplication
preserves both inequalities: \(w_i m\leq w_i y_i\leq w_i M\). Summing
over indices gives

\[
m\sum_i w_i\;\leq\;\sum_i w_i y_i\;\leq\;M\sum_i w_i.
\]

Normalization, \(\sum_i w_i=1\), gives the claimed interval. This
identifies where each assumption enters: nonnegativity preserves order,
while normalization returns the bounds to the original scale.

For unnormalized nonnegative weights with positive total \(W\), the same
argument applies to \(w_i/W\), and hence to \((\sum_i w_i y_i)/W\). If
\(W=0\), that expression is undefined. Dropping the denominator is not
an equivalent estimator and does not inherit the same range bound.

\section{Counterexamples to Broader
Claims}\label{counterexamples-to-broader-claims}

Normalization alone is insufficient. Take \((y_1,y_2)=(0,1)\) and
\((w_1,w_2)=(-1,2)\). The weights sum to one, but the aggregate is 2,
outside the input interval \([0,1]\). Negative weights invalidate the
order argument in the proof. A manuscript cannot omit nonnegativity
while retaining the proposition unchanged.

Nor does the range bound imply accuracy. Suppose two predictors both
output 100 when the true value is 0. Every convex combination of those
outputs remains 100. The proposition holds exactly, yet the squared
error is 10000. There is no contradiction: containment and accuracy are
different properties.

The example also gives no guarantee of unbiasedness. Bias requires a
target and an expectation under a specified distribution, neither of
which appears in the proposition. A paper claiming such a guarantee
would need additional assumptions and a separate argument.

\section{Verification and Limits}\label{verification-and-limits}

The repository tests finite rational examples using exact arithmetic and
checks the counterexamples numerically. Those tests are useful for
catching transcription mistakes; they do not prove the proposition for
all real inputs. The proof above supplies that argument. This
distinction is the theoretical analogue of separating a measured result
from a universal empirical claim.

This note contains no new aggregation method, empirical benchmark,
formal proof-assistant verification, or peer-review outcome. It was
authored with AI assistance as an open writing example. Its scientific
scope is the stated elementary inequality and the demonstrated limits of
its interpretation.

\section{Conclusion}\label{conclusion}

The defensible headline is that nonnegative normalized aggregation
preserves the input range. Reliability, accuracy, and generalization are
not consequences of that statement alone. A theoretical manuscript
should keep the assumptions that make its claim true visible in the
abstract, proposition, and interpretation.

\bibliography{references.bib}

\end{document}
